\chapter{Running an evaluation, and what a mask costs}\label{ch:runeval}

\begin{onething}
An evaluation here is one JSON config - models, datasets, scorers, measures - run by one command that loops scene by
scene, caches every score map under a key built from everything that could change it, and reports objects caught
beside their chance level; its first real run showed that every benchmark needs a model-free baseline in it, and its
second showed that excluding a surface class is a scope decision with a price you can count in objects.
\end{onething}

\section{Why one evaluator}\label{sec:runeval:why}

\Cref{ch:measures,ch:benchmarks} gave you the measures and the implant bench. Until 2026-10-17 each was run by its own
script (\file{scripts/evaluate_prisma_bench.py}, \file{scripts/evaluate_implant_bench.py}, code in
\file{app/utils/evaluation/}), and each read bands, reconstructed and counted slightly differently. Two scripts that
disagree in the third decimal cannot tell you whether a model improved in the third decimal. On 2026-10-17 they were
replaced by one subsystem (commit f6a444b), designed in \file{docs/lld/evaluation.md}:

\begin{center}
\small
\begin{tabular}{p{4.6cm}p{8.0cm}}
\toprule
file & what it holds \\
\midrule
\file{app/evaluation/config.py} & \code{EvaluationConfig} and its four kinds of spec (model, dataset, scorer, metric) \\
\file{app/evaluation/datasets.py} & \code{PrismaBench}, \code{ImplantBench}, \code{EvalScene} \\
\file{app/evaluation/models.py} & \code{EvalModel}: reading, reconstruction, segmentation \\
\file{app/evaluation/scorers.py} & \code{score()}: every scorer kind \\
\file{app/evaluation/metrics.py} & the measures of \cref{ch:measures}, exact, and the analytic chance \\
\file{app/evaluation/runner.py}, \file{report.py} & the loop, the cache, the tables and \file{report.md} \\
\file{scripts/evaluate.py} & the command \\
\file{configs/eval/stage3.json} & the reference config \\
\bottomrule
\end{tabular}
\end{center}

\begin{analogy}
A driving examiner builds a course of cones. A go-kart, low and twitchy, weaves through it faster than a rally driver in
a real car; on a real road the rally driver wins by a mile. Nothing was wrong with the stopwatch: the course rewarded
what go-karts are good at. You only find that out if you put a go-kart on the course on purpose. In this chapter the
go-kart is a model-free detector, PCA-RX, and the course is bench-v1.
\end{analogy}

\section{One config}\label{sec:runeval:config}

The command is
\begin{lstlisting}[language={}]
python scripts/evaluate.py configs/eval/stage3.json [--output <dir>]
\end{lstlisting}
and the config has five blocks (shortened from \file{configs/eval/stage3.json}):
\begin{lstlisting}[basicstyle=\ttfamily\footnotesize,language={}]
{"name": "stage3: ...",
 "models": [
   {"name": "v0.2.0", "checkpoint": "C:/es/models/...epoch200.pt"},
   {"name": "FT03_ep10", "checkpoint": "s3://.../FT03-joint/checkpoints/best.pt"}, ...],
 "datasets": [
   {"kind": "prisma_bench", "path": "C:/Users/FITT/Documents/stage1_datasets"},
   {"kind": "implant", "path": "C:/es/implant/bench-v1"}],
 "scorers": [
   {"kind": "residual_mahalanobis"},
   {"kind": "combined"},
   {"kind": "exclude", "classes": ["water", "cloud"],
    "segmenter": "FT03_ep10", "fit_on_kept": false,
    "base": {"kind": "residual_mahalanobis"}}, ...,
   {"kind": "pca_rx"}],
 "metrics": [
   {"kind": "roc_auc"},
   {"kind": "pauc", "max_fpr": 0.01},
   {"kind": "tpr_at", "fpr": 0.001},
   {"kind": "objects_caught", "top_pct": [1.0, 0.1]}],
 "group_by": ["tier", "source", "sensor"],
 "output": "C:/es/eval/stage3"}
\end{lstlisting}

\index{evaluation config}Four things in it need explaining.

\paragraph{A model's reading.} \index{reading (of a scene)}Each model reads a scene its own way (\code{reading}, default \code{auto}): the
wavelength-reading models read the scene's own bands within their trained-wavelength record; the shipped v0.2.0 reads
the 165-band common grid (\cref{ch:bandpipeline}). PRISMA scenes reload through the product's band filter; implant
hosts are linearly resampled, ``close to, not identical with'' the product.

\begin{sourcenote}[: the resample and the sort bug]
On PRISMA hosts the resample was far from close until commit b94dbfd: PRISMA lists its bands descending and
\code{np.interp} was called unsorted, so every grid band v0.2.0 read was a copy of one of two near-infrared bands.
Invariant 2 of \file{docs/lld/evaluation.md} now requires \code{resample} to sort first, and the fixed generic path
reproduces the product path on the real bench exactly. Every v0.2.0 implant number before 2026-10-18 is void
(\cref{ch:sortbug}).
\end{sourcenote}

\paragraph{Two new scorers.} \index{exclude scorer}\code{exclude} takes any base scorer, runs it, then sets chosen classes to 0 using a
segmentation head's class map (\file{app/evaluation/scorers.py}, lines 60--63):
\begin{lstlisting}[basicstyle=\ttfamily\footnotesize,literate={~}{{\textasciitilde}}1]
out = np.isin(classes, [CLASS_IDS[c] for c in spec.classes])   # (H, W) bool: excluded
s = score(spec.base, cube, recon, valid & ~out if spec.fit_on_kept else valid,
          classes, log_b)                                       # (H, W) base scores
return np.where(out, 0.0, s).astype(np.float32)                # excluded -> 0, still valid
\end{lstlisting}
\index{PCA-RX}\code{pca_rx} is the classical baseline: residual Mahalanobis (\cref{ch:residualmaha}) handed a reconstruction of all
zeros, so the ``residual'' is the cube itself. With the defaults \code{d = 30}, \code{local_window = (21, 61)} that is
local RX (\cref{ch:localrx}) on the scene's top 30 whitened principal components: no model, no training.

\paragraph{Labels.} Each scorer gets a name from its settings, for example
\begin{lstlisting}[language={}]
residual_mahalanobis(d=30,w=21x61)
exclude(water+cloud;mask)[residual_mahalanobis(d=30,w=21x61)]
\end{lstlisting}
\texttt{mask} when \code{fit_on_kept} is false (score the whole scene, zero the classes afterwards), \texttt{fit} when
the base scorer's statistics are fitted on the kept pixels only.

\paragraph{Validation.} The config refuses to start if two models share a name, two scorers share a label, an
\code{exclude} lacks its classes or segmenter, or a segmenter is not one of the configured models
(\code{test_config_validation}).

\section{The scene-major loop}\label{sec:runeval:loop}

\index{scene-major loop}\file{app/evaluation/runner.py} goes \textbf{scene first, then model}:
\begin{lstlisting}[basicstyle=\ttfamily\footnotesize,language={}]
for each scene (one in memory):
  for each model:  reconstruct once -> every model-based scorer -> free the reconstruction
  model-free scorers (pca_rx) once, on the scene's own cube
  class maps (for exclude) once per scene per segmenter, cached
  each score map -> the measures -> per-object rows + per-scene rows
\end{lstlisting}

Why that order: a reconstruction is as large as the cube. One PRISMA scene of about $1{,}000 \times 1{,}000$ pixels
and 185 bands\tagA{} in 4-byte floats is
\begin{align*}
185 \times 1{,}000 \times 1{,}000 &= 185{,}000{,}000 \text{ numbers}\\
185{,}000{,}000 \times 4 \text{ bytes} &= 740{,}000{,}000 \text{ bytes} \approx 0.74 \text{ GB.}
\end{align*}
The old scripts held several reconstructions at once and ran out of memory (LLD, ``Decisions''). Now memory holds one
scene and one reconstruction.

\begin{worked}[label={we:runeval:count}]{Counting the work of stage 3}
The stage-3 config has 4 models, 4 model-based scorers (residual Mahalanobis, combined, two exclusions), 1
model-free scorer, and 4 PRISMA scenes plus 18 implant hosts.
\begin{align*}
\text{score maps per scene} &= 4 \text{ models} \times 4 \text{ scorers} + 1 = 16 + 1 = 17\\
\text{scenes} &= 4 + 18 = 22\\
\text{score maps in all} &= 22 \times 17 = 374\\
\text{reconstructions} &= 22 \times 4 = 88 \quad \text{(one per scene per model)}
\end{align*}
\code{test_end_to_end_and_cache} pins the ``never more'': a stand-in model counts its calls, and after a run with
three model-based scorers it has been called once per scene.
\end{worked}

\section{The cache and its key}\label{sec:runeval:cache}

\index{score-map cache}Every score map is saved, compressed, under \file{<output>/cache/scores/}. As 16-bit floats a
$1{,}000 \times 1{,}000$ map is $1{,}000{,}000 \times 2 = 2{,}000{,}000$ bytes, about 2 MB before compression, against
740 MB for the reconstruction it came from. That ratio is why score maps are cached and reconstructions are not.

\paragraph{Float16, and its range.}\index{float16} A 16-bit float stores numbers up to 65,504; anything larger becomes infinity.
Residual Mahalanobis is a squared distance with a heavy tail, and the blurred and mixture variants of 2026-10-19
reached $10^5$. A cached infinity would put that pixel at the top of every ranking for ever, so since then a map is
stored as float16 only if all its values are finite and below $6 \times 10^4$, otherwise as float32
(\file{runner.py}, lines 115--118).

\paragraph{The key.}\index{cache key} The score map's path is
\begin{equation}\label{eq:runeval:path}
\texttt{scores/}\langle\text{dataset}\rangle\texttt{/}\langle\text{scene}\rangle\texttt{/}\langle\text{model id}\rangle\texttt{/}
\mathrm{sha1}(\text{scorer spec JSON} \,\|\, \texttt{"|"} \,\|\, \text{scene fingerprint})_{[:10]}\texttt{.npz},
\end{equation}
where $\|$ means ``joined to'' and $_{[:10]}$ the first ten hexadecimal characters (\code{_score_path}, line 144). The
\emph{model id} is the model's name followed by the first ten characters of a SHA-1 of the checkpoint's resolved path,
size, modification time and reading mode, plus adaptation settings and cloud rule when those are set
(\code{_model_id}, line 130). The \emph{scene fingerprint} is the PRISMA \texttt{.he5} file's size and modification time,
or an implant host's \texttt{cube.npy} and \texttt{gt.npy} SHA-256s from \file{manifest.json}. A
\index{hash}\emph{hash} such as SHA-1 turns any text into a fixed-length code that changes completely when any
character of the text changes, so equal keys mean equal inputs.

Change any part and the map is recomputed. Add a new measure or a new \code{group_by} and \textbf{no model runs at all}:
the measures are recomputed from cached maps.

\begin{sourcenote}[: what the key did and did not cover]
When the course part for this chapter was written, the scene's content was not in the key: a benchmark rebuilt in
place under the same scene names, evaluated into the same output directory, reused the old score maps against the new
ground truth with no error. Commit 229d423 (2026-10-18) added the fingerprint for score maps. \emph{Class maps} for
\code{exclude} are still keyed by scene name and segmenter only (\code{classes_for}, \file{runner.py} lines 80--90), so
a host rebuilt in place keeps a class map computed on its old pixels. And because the key includes the scorer spec's
full JSON, adding any field to \code{ScorerSpec} changes every key: the scoring-gains maps of 2026-10-19 could not be
reused for the bootstrap of \cref{ch:bootstrap} after \code{scale_power} and \code{scale_floor_pct} were added
(\file{docs/lld/evaluation.md}, ``Decisions'').
\end{sourcenote}

\section{Chance, per object, and the report}\label{sec:runeval:chance}

\Cref{ch:measures} derived the chance level of ``objects caught'': an object with $m$ cells after the 1-pixel dilation
is caught by random scores at budget $q$\% with probability $1 - (1 - q/100)^m$. The evaluator stores this for every
object (\code{chance_1}, \code{chance_0.1} in \file{per_object.csv}) and the report averages it into a \code{*chance*}
row\index{chance row} at the top of every table: \textbf{10.2\%} at the top 1\% on the PRISMA bench and \textbf{10.0\%} on bench-v1
\hfil\penalty50\hfilneg\tagM{docs/training/anomaly-detection/19-running-an-evaluation.md}.

\begin{worked}[label={we:runeval:chance}]{A bench's chance row}
A toy bench\tagA{} has 6 one-pixel objects ($m = 5$), 3 dominoes ($m = 8$) and 1 three-by-three square ($m = 21$). At
the top 1\%:
\begin{align*}
\text{one-pixel} &: 1 - 0.99^5 = 0.0490 && \times 6 = 0.2941\\
\text{domino} &: 1 - 0.99^8 = 0.0773 && \times 3 = 0.2318\\
\text{square} &: 1 - 0.99^{21} = 0.1903 && \times 1 = 0.1903\\
\text{chance row} &= (0.2941 + 0.2318 + 0.1903) / 10 = 0.7162 / 10 = 0.0716.
\end{align*}
About 7.2\%. A model catching 2 of these 10 objects (20\%) is about 2.8 times chance.
\end{worked}

\file{report.py} writes one section per dataset, \textbf{never pooled}\index{dataset pooling}, because their objects are different populations
(LLD, invariant 5). In each: a main table with one row per \texttt{model | scorer} and the chance row first; ``caught top
$q$\%'' is the share of all objects in that dataset; pixel measures are \emph{means over scenes}; one table per
\code{group_by} attribute. Model \texttt{-} marks a model-free scorer. The run also writes \file{per_object.csv},
\file{per_scene.csv}, \file{summary.json} and \file{provenance.json} (commit, checkpoint SHA-256s).

\section{Stage 3: the numbers held}\label{sec:runeval:stage3}

A new framework that disagrees with the old one tells you nothing until you know which is wrong. So the reference run's
first job was to reproduce every earlier figure. It did: on the PRISMA bench v0.2.0 scored pAUC 0.159, ROC-AUC 0.854
and 402 of 909 objects at the top 1\%; F001 0.148 / 0.861 / 400; FT03 0.158 / 0.853 / 413
\hfil\penalty50\hfilneg\tagM{docs/lld/evaluation.md}.

\section{The go-kart: PCA-RX wins on implants and loses on real objects}\label{sec:runeval:pcarx}

\begin{drift}
Stage 3, \file{configs/eval/stage3.json}, 2026-10-17\hfil\penalty50\hfilneg\tagM{research/experiments/implant-bench/README.md}.
\begin{center}
\small
\begin{tabular}{lccccc}
\toprule
 & \multicolumn{2}{c}{implant bench-v1} & \multicolumn{3}{c}{real PRISMA bench} \\
 & caught top 1\% & pAUC & pAUC & TPR at 0.1\% & caught top 0.1\% \\
\midrule
FT03 + res.\ Mahal. & 26.9\% & 0.172 & 0.158 & 0.058 & 14.3\% \\
PCA-RX, no model & \textbf{30.2\%} & \textbf{0.209} & 0.128 & 0.029 & 8.1\% \\
\bottomrule
\end{tabular}
\end{center}
\end{drift}

Count it in objects. On the implants PCA-RX catches $0.302 \times 3{,}336 = 1{,}007$ objects to FT03's $0.269 \times
3{,}336 = 897$. On the real objects it reaches half of FT03's detection rate at 0.1\% false alarms ($0.029 / 0.058 =
0.5$) and catches $0.081 \times 909 = 74$ objects at the top 0.1\% to FT03's $0.143 \times 909 = 130$. With bench-v1
as your only benchmark you would conclude that the models add nothing a 1990s detector \citep{reed1990adaptive} cannot
do. The real objects say the opposite.

\coursefig[0.9\linewidth]{ad-19-two-benches}{pAUC to 1\% false alarms. On implant bench-v1 (a) the model-free PCA-RX
(teal) is highest; on the real PRISMA bench (b) it is lowest, and excluding water and cloud (hatched) helps the model
on both.}{fig:runeval:twobenches}

\paragraph{Why.} \Cref{ch:benchmarks} measured it: bench-v1 implants are seven times as foreign as real objects (foreign
fraction 0.042 against 0.006), and whitening multiplies exactly foreignness. Work the multiplication once. Suppose the
scene varies with standard deviation 0.05 along its main direction and 0.002 along a minor one\tagA{}, and a pixel
differs by 0.01 along each in turn. Whitening divides each component by its direction's standard deviation and squares:
\begin{align*}
\text{main direction} &: (0.01 / 0.05)^2 = 0.2^2 = 0.04\\
\text{minor direction} &: (0.01 / 0.002)^2 = 5^2 = 25\\
\text{ratio} &: 25 / 0.04 = 625.
\end{align*}
The same size of difference scores 625 times higher when it points where the scene does not vary. PCA-RX needs nothing
more to light up a foreign implant. A model earns its keep on real objects, which are unusual mixtures of the scene's
own materials, by having learnt what each place \emph{should} look like from its surroundings.

\begin{keypoint}
Put a model-free baseline in every benchmark: if it wins there but loses on real labels, the benchmark is wrong.
\end{keypoint}

That is why bench-v2 had to pass the ordering test (\cref{sec:benchmarks:accept}) before any model was ranked or
trained on implants.

\section{Exclusion, and why ROC-AUC falls while the top improves}\label{sec:runeval:exclude}

The first product rule (user, 2026-10-17) was \textbf{``we don't detect anything on water and cloud''}. FT03's
segmentation head marks them; \code{exclude} zeroes them. \index{invariant 7}Invariant 7 of the LLD spells out the bookkeeping:
\textbf{excluded pixels stay in the measures}, as negatives tied at 0 at the bottom of the ranking, and any labelled
object inside them stays a positive that can no longer be found - ``the product would not flag it either''.

Two things happen when a large area is set to 0: its background pixels all tie at the bottom, which is harmless for the
top of the list; and any target inside it is scored 0 too, which costs it. The AUC counts the second; the top-of-list
measures mostly see the removal of bright false alarms.

\begin{worked}[label={we:runeval:mini}]{Exclusion in miniature}
Two targets, A (land) and B (on water); four background pixels, one a bright water glint. ROC-AUC is the share of
(target, background) pairs the target wins, a tie counting half (\cref{eq:measures:mw}).
\begin{align*}
&\text{before: A } 0.9,\ \text{B } 0.6;\ \text{background: glint } 0.95 \text{ (water)}, 0.5, 0.2, 0.1 \text{ (water)}\\
&\quad \text{A beats } 0.5, 0.2, 0.1 \to 3 \text{ of } 4; \quad \text{B beats } 0.5, 0.2, 0.1 \to 3 \text{ of } 4\\
&\quad \text{AUC} = (3 + 3)/8 = 0.75; \quad \text{top pixel: the glint, a false alarm}\\[4pt]
&\text{exclude water: glint} \to 0,\ 0.1 \to 0,\ \text{B} \to 0\\
&\quad \text{A beats } 0, 0.5, 0.2, 0 \to 4 \text{ of } 4\\
&\quad \text{B } (0) \text{ loses to } 0.5, 0.2; \text{ ties the two 0s} \to 0 + 0 + \tfrac12 + \tfrac12 = 1\\
&\quad \text{AUC} = (4 + 1)/8 = 0.625; \quad \text{top pixel: A, a hit.}
\end{align*}
AUC fell and the top of the list got better: the same two directions as stage 3.
\end{worked}

On the real bench, with FT03\hfil\penalty50\hfilneg\tagM{docs/lld/evaluation.md}: excluding water and cloud lifted objects caught at the top
0.1\% from 14.3\% to 16.5\% ($0.143 \times 909 = 130$ to $0.165 \times 909 = 150$, 20 more objects in the same review
budget) and pAUC from 0.158 to 0.165, while ROC-AUC fell from 0.853 to 0.830. Most of the fall is set1, which FT03
calls 55\% water (ROC-AUC 0.932 to 0.851 there). Every model gained at the top 0.1\% by about the same amount, including
v0.2.0, which has no segmentation head: \textbf{the gain belongs to the mask, not to any model}.

\section{Scope: cloud excluded, water kept}\label{sec:runeval:scope}

\begin{analogy}
A security guard is told to ignore the car park, because it is full of harmless reflections. Her shift gets quieter. But
the day somebody breaks into a car, nobody sees it - the instruction did not make the car park safe, it made it
invisible. And if her map of the car park is drawn too large, she also ignores half the loading bay.
\end{analogy}

\index{scope}So the question ``should we exclude water?'' is not a modelling question. It is: \textbf{are objects on water something
the product must find?} For vessels, yes.

\paragraph{The claim that was wrong.} Before the decision the record said ``no real object sits on water''. It came
from a scan that called a pixel water when its visible reflectance (mean over roughly 450--650 nm) was below 0.04.
Clear open water is that dark; turbid coastal water, with silt or a shallow bottom, is around 0.08.

\begin{worked}[label={we:runeval:water}]{Two water rules on one turbid pixel}
A pixel with visible 0.08, red 0.06, near-infrared 0.03\tagA{}.
\begin{align*}
\text{scan's rule: } & \mathrm{VIS} < 0.04? \quad 0.08 < 0.04 \text{ is false} \to \text{``land''}\\
\mathrm{NDVI} &= \frac{\mathrm{NIR} - \mathrm{red}}{\mathrm{NIR} + \mathrm{red}} = \frac{0.03 - 0.06}{0.03 + 0.06}
  = \frac{-0.03}{0.09} = -0.33\\
\text{builder's rule: } & \mathrm{NIR} < 0.08 \text{ and } \mathrm{NDVI} < 0? \quad \text{yes and yes} \to \text{water.}
\end{align*}
Water is dark in the near-infrared even when it is bright in the visible; the scan used the wrong band.
\end{worked}

Re-checked on 2026-10-17 against FT03's water / cloud mask, \textbf{39 of the 909 real objects lie wholly inside
it}\hfil\penalty50\hfilneg\tagM{docs/lld/evaluation.md}:

\begin{center}
\small
\begin{tabular}{p{5.8cm}cp{5.2cm}}
\toprule
where & objects & what they are \\
\midrule
set1, on water, more than 3 px offshore & 13 & single pixels; likely vessels \\
set1, on water, within 3 px of shore & 14 & single pixels near the coast \\
set2 and set3, shorelines & 8 & objects on the land--water edge \\
set4, under haze FT03 calls cloud & 4 & a scope question at the haze edge \\
\midrule
total & 39 & \\
\bottomrule
\end{tabular}
\end{center}

\coursefig[0.85\linewidth]{ad-21-the-39}{The 39 labelled objects lying wholly inside FT03's water / cloud mask, by kind.
Excluding water and cloud hides all 39 (amber); excluding cloud only hides the 4 in set4 that sit under haze FT03 calls
cloud.}{fig:runeval:the39}

The user's rule then changed: \textbf{vessels on water are in scope; only cloud is excluded.}

\paragraph{What each choice buys, counted in objects.} FT03, residual Mahalanobis, PRISMA bench
(\file{C:/es/eval/cloud\_only/report.md} on the dev laptop):

\begin{drift}
\file{configs/eval/cloud_only.json}, 2026-10-17\hfil\penalty50\hfilneg\tagM{docs/training/anomaly-detection/21-scope-and-masks.md}.
Objects caught at the top 0.1\%, per scene, from \file{per_object.csv}:
\begin{center}
\small
\begin{tabular}{lccccc}
\toprule
 & set1 & set2 & set3 & set4 & total \\
\midrule
plain & 56 & 33 & 7 & 34 & 130 \\
cloud only & 56 & 33 & 13 & 42 & 144 \\
water + cloud & 59 & 34 & 14 & 43 & 150 \\
\bottomrule
\end{tabular}
\end{center}
pAUC 0.158 / 0.166 / 0.165 and ROC-AUC 0.853 / 0.850 / 0.830 (plain / cloud only / water + cloud).
\end{drift}

Work the shares: $130/909 = 14.3\%$, $144/909 = 15.8\%$, $150/909 = 16.5\%$. Cloud exclusion buys 14 objects
($144 - 130$: set3 $+6$, set4 $+8$) and costs almost no ROC-AUC. Adding water buys 6 more at the top 0.1\% but none at
the top 1\% (both catch $419/909 = 46.1\%$), takes ROC-AUC from 0.850 to 0.830, and hides the vessels. pAUC, the
measure the analyst's view tracks, is no better for it (0.165 against 0.166). \textbf{Water exclusion trades the
vessels for a slightly tidier top of the list}; with vessels in scope, that is a bad trade. \Cref{ch:bootstrap} shows
that cloud exclusion's gain is one of the few on this bench whose interval clears zero.

\paragraph{The set4 haze, and a claim that was wrong twice.}\index{haze colour test} Cloud-only exclusion still hides 4 objects in set4. On
2026-10-17 that was read as a mask error, ``FT03 calls bright desert cloud''. A look at the true-colour picture on
2026-10-18 showed a bright cloud/smoke streak and a large hazy patch that FT03's 11\% follows. The colour test settles
it: bright desert is red, its reflectance rising with wavelength, so near-infrared over blue is above 1; haze scatters
blue most, so near-infrared over blue falls below 1.
\begin{align*}
\text{desert-like: } & \text{blue } 0.20, \text{ NIR } 0.35 \to 0.35 / 0.20 = 1.75 \quad \text{(red)}\\
\text{set4 ``cloud'': } & \text{blue } 0.30, \text{ NIR } 0.18 \to 0.18 / 0.30 = 0.60 \quad \text{(bluish)}
\end{align*}
For 75\% of set4's ``cloud'' pixels the ratio is 0.54--0.65
\hfil\penalty50\hfilneg\tagM{docs/training/anomaly-detection/21-scope-and-masks.md}: haze. The 4 objects are things a labeller could see
through haze, a scope question (should the product look under thin cloud?), not a mask fault. A negative claim about a
picture deserves the same check as a positive one: look at the picture.

\paragraph{The arithmetic of a wrong mask.} The next mask may be wrong, so know its price. Had 11\% of set4 been called
cloud wrongly, the analyst would lose 11\% of the scene, holding 4 of set4's 230 objects
\hfil\penalty50\hfilneg\tagM{docs/training/anomaly-detection/21-scope-and-masks.md} ($4 / 230 = 1.7\%$). Small
here; but on a scene full of targets a wrong mask loses them in proportion, \emph{silently} - the score is 0, not
``unknown''.

\begin{keypoint}
Decide scope first, then count the labelled objects inside the mask, then audit the mask against the picture, and only
then judge the result at the top of the list.
\end{keypoint}

\begin{confusions}
\confusion{PCA-RX won on the implants, so the models are no better than classical detection?}{It won on implants built
in a way that suits it. On the 909 real objects it is last. The benchmark was wrong; the stopwatch was fine.}
\confusion{Exclusion lowered ROC-AUC, so it made detection worse?}{It tied a large area at 0 and hid the objects inside
it; the top of the list, where the analyst looks, improved for every model. Judge it by pAUC, TPR at 0.1\% and objects
caught, with ROC-AUC beside them.}
\confusion{I added a measure and re-ran; the numbers cannot be stale?}{Score maps are keyed by scene fingerprint since
229d423; class maps are not. Give each benchmark version its own path and output directory.}
\confusion{\texttt{exclude} with segmenter FT03 only applies to FT03?}{It applies to every model in the config; the
segmenter only supplies the class map.}
\confusion{The PRISMA bench and the implant bench are averaged in the summary?}{Never. Each dataset gets its own
section and its own chance row.}
\confusion{Cloud-only exclusion also hides the 39 objects?}{No: 35 of them are inside the water mask. Cloud-only hides
the 4 in set4 under haze.}
\end{confusions}

\begin{checkyourself}
\item\label{ex:runeval:1} A config has 3 models, 3 model-based scorers and 1 model-free scorer, over 5 scenes. How many
score maps and reconstructions does a first run make? And a second run after adding one measure?
\item\label{ex:runeval:2} In \cref{we:runeval:mini}, suppose B had been on land (not excluded) at 0.6. What are the AUC
and the top pixel after excluding water?
\item\label{ex:runeval:3} A pixel has visible 0.07, red 0.05, NIR 0.02. What do the scan's rule and the builder's rule
call it?
\item\label{ex:runeval:4} At the top 1\%, plain FT03 catches 101 / 103 / 48 / 161 objects on set1--4, cloud-only 101 /
103 / 54 / 161, water + cloud 98 / 107 / 53 / 161. Compute each total and share of 909. What does water exclusion swap?
\item\label{ex:runeval:5} A score map's largest value is 70,000. What happens if it is cached as float16, and what does
the runner do instead? How many bytes does a $1{,}000 \times 1{,}000$ map take uncompressed in each format?
\item\label{ex:runeval:6} Put the four questions of the last key point in order for a proposal to exclude snow.
\end{checkyourself}

\nextchapter{ch:bootstrap}
